% ── Appendix A — The Snowflake Schema ────────────────────────────────────────
\chapter{The Snowflake Schema: A Normalized Dimensional Variant}
\label{app:snowflake}

Chapter~1 presented the Kimball star schema (Figure~\ref{fig:ch1-star}), in which a central fact table is surrounded by \emph{denormalized} dimension tables: each dimension holds all of its descriptive attributes in a single flat table. The \emph{snowflake schema} is the normalized variant of that design. Each dimension is decomposed into a hierarchy of related tables, so that repeating attribute groups are stored once and referenced by foreign keys, exactly as third normal form prescribes \parencite{kimball2013}. This appendix presents the snowflake schema as a reference, because it sits between the two poles the thesis compares --- the query-optimized denormalization of Kimball and the change-optimized normalization of Data Vault~2.0 --- and so illustrates the structural-scalability trade-off concretely.

\fig[0.98\textwidth]{snowflake_schema.png}{A snowflake schema for a retail sales process. The central \texttt{Fact\_Sales} table references the \texttt{Dim\_Date}, \texttt{Dim\_Store}, and \texttt{Dim\_Product} dimensions, each of which is further normalized into sub-dimensions (for example \texttt{Dim\_Month} and \texttt{Dim\_Quarter} off \texttt{Dim\_Date}, and \texttt{Dim\_Brand} off \texttt{Dim\_Product}). \textit{Source: SqlPac, Wikimedia Commons (CC BY-SA 3.0).}}{app-snowflake}

Figure~\ref{fig:app-snowflake} shows the same retail sales process modeled as a snowflake: where a star schema would store the product's brand and category as columns inside one \texttt{Dim\_Product} table, the snowflake factors them into separate \texttt{Dim\_Brand} and \texttt{Dim\_Product\_Category} tables linked by foreign keys.

\section*{How Snowflaking Changes the Trade-off}
\addcontentsline{toc}{section}{How Snowflaking Changes the Trade-off}

Normalizing the dimensions buys two structural advantages and imposes one query cost. The advantages are storage efficiency --- a brand name is stored once rather than repeated on every product row --- and update integrity, since a corrected brand name propagates from a single place rather than requiring a sweep across a wide dimension table. The cost is query complexity: reconstructing a business view now requires joining the fact table to a dimension and then to its sub-dimensions, so a snowflake query carries more joins than the equivalent star query. This is the same tension, in miniature, that Chapter~1 traced between Kimball and Data Vault~2.0: normalization improves the economics of change and storage at the expense of read-time joins, and denormalization makes the opposite trade. Table~\ref{tab:app-star-snowflake} summarizes the comparison.

\begin{table}[H]
\centering
\renewcommand{\arraystretch}{1.35}
\rowcolors{2}{white}{tblrowalt}
\caption{Star versus snowflake schema: the denormalization--normalization trade-off within dimensional modeling. \textit{Author-generated, after \parencite{kimball2013}.}}
\label{tab:app-star-snowflake}
\begin{tabularx}{\linewidth}{|>{\raggedright\arraybackslash}p{3.4cm}|>{\raggedright\arraybackslash}X|>{\raggedright\arraybackslash}X|}
\hline
\rowcolor{tblheader}
\textcolor{white}{\textbf{Property}} &
\textcolor{white}{\textbf{Star schema}} &
\textcolor{white}{\textbf{Snowflake schema}} \\
\hline
Dimension form & Denormalized (one flat table) & Normalized into sub-dimensions \\
Storage of repeated values & Repeated on every row & Stored once, referenced by key \\
Query joins & Fewer; fact to dimension only & More; fact to dimension to sub-dimension \\
Read performance & Faster for interactive BI & Slower; more joins per query \\
Update integrity & Wide updates across the dimension & Localized to one small table \\
\hline
\end{tabularx}
\end{table}

\section*{Relevance to the Thesis}
\addcontentsline{toc}{section}{Relevance to the Thesis}

The snowflake schema does not change the verdict of Chapter~1. Both star and snowflake remain dimensional designs built for a stable, pre-specified analytical perimeter, and both inherit the schema rigidity that the chapter identified: a structural change to a widely shared dimension still propagates to every fact and report that uses it. Snowflaking only redistributes attributes across more tables; it does not make structural change \emph{additive} in the sense that the Hub--Link--Satellite decomposition does. The schema is included here as a reference point that makes the normalization axis explicit: the same principle that Data Vault~2.0 carries to its conclusion --- store each thing once, in its own object --- appears in milder form in the snowflake, which is why practitioners sometimes describe Data Vault as occupying the far end of the normalization spectrum the star and snowflake begin.

\medskip
\noindent\footnotesize The figure is reproduced from Wikimedia Commons under the Creative Commons Attribution-ShareAlike 3.0 licence (CC BY-SA 3.0); author: SqlPac. Source: \url{https://commons.wikimedia.org/wiki/File:Snowflake-schema-example.png}.
